\documentclass[main.tex]{subfiles}
\begin{document}
% ============================================================
%  Appendix: adjoint calculus
% ============================================================
\section{Appendix: Adjoint Calculus}

\begin{frame}[label=adjoints]{Adjoint Calculus: Sensitivities vs.\ Adjoints}
  \begin{hlbox}
    After discretization $z \in \R^m$ with $m$ large. Two ways to compute $\hat J'(z)$:
  \end{hlbox}

  \begin{columns}[T]
    \begin{column}{0.48\textwidth}
      \begin{hlbox}
        \textbf{Sensitivities.} For each direction $h_i$, solve $e_u\,\delta u_i = -e_z h_i$, then
        \[
          \hat J'(z)h_i = J_u\,\delta u_i + J_z h_i .
        \]
        \textbf{Cost:} $m$ linearized solves.
      \end{hlbox}
    \end{column}
    \begin{column}{0.48\textwidth}
      \begin{hlbox}
        \textbf{Adjoints.} Solve once $e_u^* p = J_u$, then
        \[
          \hat J'(z) = J_z - e_z^* p .
        \]
        \textbf{Cost:} one adjoint solve, independent of $m$.
      \end{hlbox}
    \end{column}
  \end{columns}

  \begin{qbox}
    Sensitivities pay off only for few controls, or for many outputs (e.g., many constraints).
  \end{qbox}
  \vfill\hfill\hyperlink{hessian<3>}{\beamerreturnbutton{Back to Section 5}}
\end{frame}

\begin{frame}{Adjoint Calculus: Second Derivatives}
  \begin{hlbox}
    \textbf{Where the recipe comes from.} For $J$, $e$ of class $C^2$, with $p = p(z)$ the adjoint state and $w(h) := (S'(z)h,\, h)$,
    \[
      \hat J''(z)[h_1,h_2] = \Lag_{(u,z)(u,z)}(S(z),z,p)\big[w(h_1),\, w(h_2)\big] .
    \]
  \end{hlbox}

  \begin{hlbox}
    \textbf{Idea.} Differentiate $\hat J(z) = \Lag(S(z),z,p(z))$ twice. The terms with $S''(z)$ and $p'(z)$ drop out because $\Lag_u = 0$ (adjoint equation) and $\Lag_p = 0$ (state equation).
  \end{hlbox}

  \begin{hlbox}
    \begin{itemize}
      \item Second derivatives of $e$ enter only through $\dual{p}{e''(u,z)[\cdot,\cdot]}$.
      \item The Hessian is dense and $m \times m$: never form it. Newton--Krylov methods only need $\hat J''(z)h$.
    \end{itemize}
  \end{hlbox}
  \vfill\hfill\hyperlink{hessian<3>}{\beamerreturnbutton{Back to Section 5}}
\end{frame}

\begin{frame}{Adjoint Calculus: Second-Order Sufficient Conditions}
  \begin{hlbox}
    \textbf{Convex case} (Examples 1, 2): first-order conditions are already sufficient.
  \end{hlbox}

  \begin{hlbox}
    \textbf{Nonconvex case} (Example 3): for box constraints, require $\hat J''(\bar z)[h,h] \ge \delta\norm{h}_{\Ltwo}^2$, $\delta > 0$, on the critical cone
    \[
      C(\bar z) = \{ h : h \ge 0 \text{ where } \bar z = z_a,\; h \le 0 \text{ where } \bar z = z_b,\; h = 0 \text{ where } \nabla\hat J(\bar z) \ne 0 \}.
    \]
  \end{hlbox}

  \begin{qbox}
    Subtlety: the \textbf{two-norm discrepancy}. $\hat J$ is differentiable in $L^\infty$, but coercive only in $L^2$.
  \end{qbox}
  \vfill\hfill\hyperlink{hessian<3>}{\beamerreturnbutton{Back to Section 5}}
\end{frame}

\begin{frame}{Adjoint Calculus: The Time-Dependent Problem}
  \begin{hlbox}
    \textbf{Example 2.} For $J(u,z) = \tfrac12\norm{u(T) - u_T}_{\Ltwo}^2 + \tfrac{\alpha}{2}\norm{z}_{L^2(Q_T)}^2$, the adjoint runs \textbf{backward in time}:
    \[
      -\partial_t p - \Delta p = 0 \;\;\text{in } Q_T, \qquad p = 0 \;\;\text{on } \partial D\times(0,T), \qquad p(T) = u(T) - u_T .
    \]
    The gradient in $L^2(Q_T)$ is again $\nabla\hat J(z) = p + \alpha z$.
  \end{hlbox}

  \begin{hlbox}
    \begin{itemize}
      \item The adjoint needs the \emph{whole} forward trajectory $u(t)$: store it, or recompute it (checkpointing).
      \item Optimize-then-discretize and discretize-then-optimize give the same discrete gradient only for compatible time-stepping schemes.
    \end{itemize}
  \end{hlbox}
  \vfill\hfill\hyperlink{hessian<3>}{\beamerreturnbutton{Back to Section 5}}
\end{frame}

\end{document}
